\chapter{Non-Functional Requirements}
\label{ch:non-functional-requirements}

Non-functional requirements are stated only where the repository provides
evidence, or are explicitly marked as RECOMMENDED.

\section{Compatibility}

\begin{longtable}{@{}L{4.2cm}L{10cm}@{}}
\caption{Compatibility requirements}\\
\toprule
\textbf{Area} & \textbf{Requirement / evidence} \\
\midrule
\endfirsthead
\toprule
\textbf{Area} & \textbf{Requirement / evidence} \\
\midrule
\endhead
Node.js & $\geq$ 18 (\texttt{engines.node} in \texttt{@expo-snips/cli}; esbuild
target \texttt{node18}). \\
Language target & ES2020, CommonJS modules (\texttt{tsconfig.base.json}). \\
VS Code & $\geq$ 1.80.0 (\texttt{engines.vscode}). \\
Operating system & No OS-specific code; paths are built with \texttt{path.join}
and the CLI uses \texttt{process.cwd()}, so Windows, macOS and Linux are supported. \\
Target projects & Expo / React Native with Expo Router (default); NativeWind or
StyleSheet; TypeScript or JavaScript. \\
\bottomrule
\end{longtable}

\section{Performance}

\begin{itemize}
  \item A remote registry request is bounded by a 10-second abort timeout
        (\texttt{registry.ts} in both clients).
  \item The extension caches a successful remote index for 5 minutes and enters a
        5-minute back-off after an unreachable registry, preventing repeated stalls
        (\texttt{packages/vscode-extension/src/registry.ts}).
  \item Rendering is in-memory Handlebars compilation per file; there is no
        incremental compilation or caching. This is INFERRED as adequate because
        recipes contain at most seven file entries, but no benchmark exists in the
        repository.
  \item The diff algorithm is a full $O(n \times m)$ longest-common-subsequence
        table (\texttt{diff.ts}), allocated per file compared.
\end{itemize}

Quantified performance targets (\textbf{UNKNOWN}) --- the repository defines none.

\section{Reliability and fault tolerance}

\begin{itemize}
  \item Catalog resolution degrades gracefully: remote failure falls back to the
        bundled catalog (FR-1).
  \item An unreadable or corrupt manifest is treated as empty
        (\texttt{manifest.ts}, lines 37--45).
  \item Update skips recipes that are no longer present in the registry with a
        warning instead of failing the whole operation (\texttt{update.ts},
        lines 78--83).
  \item An unknown primitive id logs a warning and is skipped
        (\texttt{primitives.ts}, line 31).
\end{itemize}

\textbf{Status:} IMPLEMENTED.

\section{Security (summary)}

Detailed in Chapter~\ref{ch:security}. Non-functional expectations:

\begin{itemize}
  \item No credentials are stored or transmitted by the product.
  \item Network access is read-only HTTP GET against JSON endpoints.
  \item Recipe content and recipe dependency names are treated as trusted input;
        there is no signature verification.
\end{itemize}

\textbf{Status:} IMPLEMENTED (behaviour) / RECOMMENDED (hardening measures).

\section{Usability}

\begin{itemize}
  \item Coloured terminal output groups information, warnings and errors
        (\texttt{chalk}).
  \item Interactive prompts are used only when a value cannot be resolved.
  \item Post-generation guidance lists concrete next steps
        (\texttt{add.ts}, lines 275--278).
  \item The extension reports progress in a notification and offers a single
        follow-up action.
\end{itemize}

\textbf{Status:} IMPLEMENTED.

\section{Accessibility}

\textbf{NOT APPLICABLE} for a graphical user interface: the product has no web or
mobile UI of its own; its interface is a command-line program and a VS Code
extension surface built entirely from native VS Code components (Quick Pick,
notification, status bar), which inherit the editor's own accessibility
behaviour.

Accessibility of \emph{generated screens} is the responsibility of the recipe
authors. The templates verify: they use \texttt{Text} components for all copy and
\texttt{Pressable} for interaction, but no recipe adds
\texttt{accessibilityLabel}, \texttt{accessibilityRole} or
\texttt{accessibilityHint}. This is recorded as a recommendation in
Chapter~\ref{ch:ui-ux} and as technical debt in Chapter~\ref{ch:technical-debt}.

\section{Maintainability}

\begin{itemize}
  \item The engine, CLI, registry and extension are separated into distinct
        packages with a single dependency direction
        (CLI/extension/registry $\rightarrow$ core).
  \item The recipe format is captured in a JSON Schema and enforced by a
        validation script.
  \item TypeScript strict mode is enabled workspace-wide
        (\texttt{strict: true}).
  \item Counter-evidence: the stack-detection logic and the primitive-resolution
        logic are duplicated between the CLI and the extension; there is no test
        suite; and generated catalog modules are committed/ignored inconsistently.
\end{itemize}

\textbf{Status:} IMPLEMENTED (structure) / INFERRED (assessment).

\section{Portability and distribution}

\begin{itemize}
  \item The CLI is published to npm with public access.
  \item The extension is packaged with \texttt{vsce} into a \texttt{.vsix}.
  \item The registry is plain static files, deployable to any static host.
  \item The catalog is additionally embedded in both clients, so distribution of
        the catalog does not depend on hosting availability.
\end{itemize}

\textbf{Status:} IMPLEMENTED.

\section{Observability}

The product emits console output only. There is no logging framework, no log
levels, no log destination configuration, no metrics and no telemetry. Diagnostics
are limited to \texttt{console.log}/\texttt{console.warn}/\texttt{console.error}
with \texttt{chalk} colouring in the CLI.

\textbf{Status:} IMPLEMENTED (as a limitation); instrumentation is RECOMMENDED for
future work but is not required for a local code-generation tool.

\section{Privacy}

\begin{itemize}
  \item No personal data is collected, transmitted or stored by the product.
  \item The only persistent artefact is \texttt{.expo-snips.json} in the user's
        project, containing recipe ids, versions, timestamps and option values.
  \item No analytics or third-party trackers exist in the repository.
\end{itemize}

\textbf{Status:} IMPLEMENTED (verified by absence).

\section{Localisation}

\textbf{NOT APPLICABLE.} All user-facing strings are hard-coded English literals;
there is no i18n framework, message catalogue or locale handling.

\section{Portability of generated code}

Generated screens target React Native only. They import from \texttt{react-native},
\texttt{expo-router} or \texttt{@react-navigation/native}, and the third-party
libraries declared by each recipe. They are therefore portable across the platforms
Expo supports, subject to the recipe's dependency set (for example
\texttt{map/location} declares \texttt{react-native-maps}, which is not available
on the web target without an additional provider).
